
\section{Conclusion}\label{sec:Conclusion}

In this work, we presented concrete Haskell implementations of two knowing-how logics: the basic logic $\mathcal{L}_{Kh}$ and its extension with regularity constraints $reg\text{-}\mathcal{L}^U_{Kh}$. For $\mathcal{L}_{Kh}$, we implemented an explicit model checker based on bounded plan search. For $reg\text{-}\mathcal{L}^U_{Kh}$, we implemented a model checker following the procedure described in \cite{Demri2023}. Since the original paper does not spell out all implementation details of the reachability component, we made these constructions explicit by using depth-first search. We also incorporated a small optimisation in the implementation of \texttt{CheckSE}: instead of performing reachability tests inside nested loops, we first compute the set of bad points and then perform the reachability checks, thereby avoiding repeated reachability computations. 

To support experimentation, we provided parsers for both logics, generators for random models and formulas, a collection of QuickCheck and HSpec tests, and a lightweight web interface for interactive use. Together, these components turn the underlying theoretical constructions into an executable and accessible system.

We note, however, that the current implementation assumes that all automata provided in the uncertainty component are well-formed. In particular, we do not explicitly verify whether they satisfy the structural conditions required by the theoretical framework. As a result, it remains the responsibility of the user to ensure that the automata used in the model conform to these requirements. A natural direction for future work would therefore be to implement a dedicated sanity checker for automata.

